% GENERATED FILE. Edit canonical lessons, then run npm run generate:books.
% canonical-source-hash: fnv1a64:4e1c03a1a0bd2268
% canonical-lessons: GE-C417-sicherheitskontrolle, GE-C417-durchsuchen, GE-C417-zoll, GE-C417-verzollen, GE-C417-reisepass

\chapter{Kontrolle und Zoll}
\label{ch:ge-kontrolle-und-zoll}
\hlchaptermodality{417}

\begin{chapteropening}
\textbf{By the end of this chapter:} \emph{I can go through security and customs: the security check, searching bags, customs, declaring goods and the passport.}

Complete the last lesson of chapter 417: I can go through security and customs: the security check, searching bags, customs, declaring goods and the passport.
\end{chapteropening}

\section[die Sicherheitskontrolle]{die Sicherheitskontrolle — the security check}

\label{lesson:GE-C417-sicherheitskontrolle}



\begin{quote}
Before the new one: say the German for the hand luggage, then the German for to give up.
\end{quote}

\subsection*{You'll want to know: die Sicherheitskontrolle}
\textbf{die Sicherheitskontrolle} — \textquotedblleft{}the security check\textquotedblright{}.

\begin{quote}
An der Sicherheitskontrolle muss er seinen Gürtel ausziehen.
\end{quote}

\begin{quote}
\emph{At the security check he has to take off his belt.}
\end{quote}

Plural: \emph{die Sicherheitskontrollen}.

\begin{cousinweb}
\emph{Sicherheitskontrolle} is built from \emph{Sicherheit} 'safety, security', from \emph{sicher}, which goes back to Latin \emph{sēcūrus} (Lewis \& Short), and \emph{Kontrolle} 'check', borrowed from French \emph{contrôle} (EtymDB).
\end{cousinweb}

\begin{grammarlens}[title={What you've built}]
One more word for telling, explaining and describing at B1.
\end{grammarlens}

\subsection*{Your turn}
\begin{itemize}
\raggedright
  \item \emph{Say it:} \emph{die Sicherheitskontrolle}
  \item \emph{Say it:} \emph{die Sicherheitskontrolle}, once more
  \item \emph{Say it:} a sentence of your own with \emph{die Sicherheitskontrolle}
  \item \emph{Recall:} say the German for the hand luggage, then the German for to give up, then say \emph{die Sicherheitskontrolle} again
\end{itemize}

\begin{tcolorbox}[breakable,colback=teal!4,colframe=teal!35!black,title={Before you move on}]
What does die Sicherheitskontrolle mean? (\textquotedblleft{}The security check\textquotedblright{}.) Say it once more.
\end{tcolorbox}
\section[durchsuchen]{durchsuchen — to search (a bag, a place)}

\label{lesson:GE-C417-durchsuchen}



\begin{quote}
Before the new one: say the German for to give up, then the German for the security check.
\end{quote}

\subsection*{You'll want to know: durchsuchen}
\textbf{durchsuchen} — \textquotedblleft{}to search (a bag, a place)\textquotedblright{}.

\begin{quote}
Eine Beamtin durchsucht ihren Rucksack.
\end{quote}

\begin{quote}
\emph{An official searches her rucksack.}
\end{quote}

Inseparable: \emph{sie durchsucht}, \emph{sie hat durchsucht}.

\begin{cousinweb}
\emph{durchsuchen} is built from \emph{durch} 'through' and \emph{suchen} 'to look for', inherited from Old High German \emph{suohhen} and Proto-Germanic \emph{sōkijaną} (EtymDB). English \emph{seek} is a cousin.
\end{cousinweb}

\begin{grammarlens}[title={What you've built}]
One more word for telling, explaining and describing at B1.
\end{grammarlens}

\subsection*{Your turn}
\begin{itemize}
\raggedright
  \item \emph{Say it:} \emph{durchsuchen}
  \item \emph{Say it:} \emph{durchsuchen}, once more
  \item \emph{Say it:} a sentence of your own with \emph{durchsuchen}
  \item \emph{Recall:} say the German for to give up, then the German for the security check, then say \emph{durchsuchen} again
\end{itemize}

\begin{tcolorbox}[breakable,colback=teal!4,colframe=teal!35!black,title={Before you move on}]
What does durchsuchen mean? (\textquotedblleft{}To search (a bag, a place)\textquotedblright{}.) Say it once more.
\end{tcolorbox}
\section[der Zoll]{der Zoll — the customs; the duty (tax)}

\label{lesson:GE-C417-zoll}



\begin{quote}
Before the new one: say the German for the security check, then the German for to search.
\end{quote}

\subsection*{You'll want to know: der Zoll}
\textbf{der Zoll} — \textquotedblleft{}the customs; the duty (tax)\textquotedblright{}.

\begin{quote}
Beim Zoll fragt ein Beamter nach Zigaretten.
\end{quote}

\begin{quote}
\emph{At customs an officer asks about cigarettes.}
\end{quote}

Plural: \emph{die Zölle}, for the taxes. \emph{Der Zoll} is also an inch, a different word.

\begin{cousinweb}
\emph{Zoll} comes from Old High German \emph{zol} (EtymDB), and it is usually traced to Late Latin \emph{tolōneum} 'toll house', from Greek \emph{tel\'{\={o}}nion}. EtymDB instead links it with a Germanic word for what is counted, so the line is hedged. English \emph{toll} is a cousin.
\end{cousinweb}

\begin{grammarlens}[title={What you've built}]
One more word for telling, explaining and describing at B1.
\end{grammarlens}

\subsection*{Your turn}
\begin{itemize}
\raggedright
  \item \emph{Say it:} \emph{der Zoll}
  \item \emph{Say it:} \emph{der Zoll}, once more
  \item \emph{Say it:} a sentence of your own with \emph{der Zoll}
  \item \emph{Recall:} say the German for the security check, then the German for to search, then say \emph{der Zoll} again
\end{itemize}

\begin{tcolorbox}[breakable,colback=teal!4,colframe=teal!35!black,title={Before you move on}]
What does der Zoll mean? (\textquotedblleft{}The customs; the duty (tax)\textquotedblright{}.) Say it once more.
\end{tcolorbox}
\section[verzollen]{verzollen — to declare (goods at customs); to pay duty on}

\label{lesson:GE-C417-verzollen}



\begin{quote}
Before the new one: say the German for to search, then the German for the customs.
\end{quote}

\subsection*{You'll want to know: verzollen}
\textbf{verzollen} — \textquotedblleft{}to declare (goods at customs); to pay duty on\textquotedblright{}.

\begin{quote}
Wir haben nichts zu verzollen.
\end{quote}

\begin{quote}
\emph{We have nothing to declare.}
\end{quote}

\emph{Nichts zu verzollen} — nothing to declare, the green exit.

\begin{cousinweb}
\emph{verzollen} is built from \emph{ver-} and \emph{Zoll} 'customs duty', from Old High German \emph{zol} (EtymDB), a word usually traced to Late Latin \emph{tolōneum}.
\end{cousinweb}

\begin{grammarlens}[title={What you've built}]
One more word for telling, explaining and describing at B1.
\end{grammarlens}

\subsection*{Your turn}
\begin{itemize}
\raggedright
  \item \emph{Say it:} \emph{verzollen}
  \item \emph{Say it:} \emph{verzollen}, once more
  \item \emph{Say it:} a sentence of your own with \emph{verzollen}
  \item \emph{Recall:} say the German for to search, then the German for the customs, then say \emph{verzollen} again
\end{itemize}

\begin{tcolorbox}[breakable,colback=teal!4,colframe=teal!35!black,title={Before you move on}]
What does verzollen mean? (\textquotedblleft{}To declare (goods at customs); to pay duty on\textquotedblright{}.) Say it once more.
\end{tcolorbox}
\section[der Reisepass]{der Reisepass — the passport}

\label{lesson:GE-C417-reisepass}



\begin{quote}
Before the new one: say the German for the customs, then the German for to declare.
\end{quote}

\subsection*{You'll want to know: der Reisepass}
\textbf{der Reisepass} — \textquotedblleft{}the passport\textquotedblright{}.

\begin{quote}
Für diese Reise braucht sie einen Reisepass.
\end{quote}

\begin{quote}
\emph{For this trip she needs a passport.}
\end{quote}

Plural: \emph{die Reisepässe}. Often shortened to \emph{der Pass}.

\begin{cousinweb}
\emph{Reisepass} is built from \emph{Reise} 'journey', inherited from Middle High German \emph{reise} (EtymDB), and \emph{Pass} 'pass, passport', which came through Italian \emph{passaporto} from Latin \emph{passus} 'step' (EtymDB).
\end{cousinweb}

\begin{grammarlens}[title={What you've built}]
One more word for telling, explaining and describing at B1.
\end{grammarlens}

\subsection*{Your turn}
\begin{itemize}
\raggedright
  \item \emph{Say it:} \emph{der Reisepass}
  \item \emph{Say it:} \emph{der Reisepass}, once more
  \item \emph{Say it:} a sentence of your own with \emph{der Reisepass}
  \item \emph{Recall:} say the German for the customs, then the German for to declare, then say \emph{der Reisepass} again
\end{itemize}

\begin{tcolorbox}[breakable,colback=teal!4,colframe=teal!35!black,title={Before you move on}]
What does der Reisepass mean? (\textquotedblleft{}The passport\textquotedblright{}.) Say it once more.
\end{tcolorbox}
